% ============================================================
% Section 3: Motivating Observations
% The Information Budget of Graph Neural Networks
% ============================================================

\section{Motivating Observations}
\label{sec:observation}

Before developing our theoretical framework (Section~\ref{sec:snr_theory_rigorous}), we present three empirical observations that reveal fundamental patterns in GNN behavior. These observations motivate the key quantities---$\kappa$, $\mathcal{B}$, and ADR---that our theory will formalize.

\subsection{Observation 1: The U-Shape Pattern in Controlled Experiments}
\label{sec:ushape_obs}

To isolate the effect of homophily, we conduct controlled synthetic experiments varying edge homophily $h \in \{0.1, 0.2, ..., 0.9\}$ while keeping features constant.

\textbf{Experimental Setup.} We generate graphs with $n = 1000$ nodes, $d = 20$ dimensional features with controlled separability, and approximately 7,500 edges achieving target homophily. Models (MLP, GCN, GAT, GraphSAGE) use identical architectures (64 hidden units, 2 layers, dropout 0.5). We report mean accuracy over 10 runs.

\begin{figure}[t]
    \centering
    \includegraphics[width=\columnwidth]{code/figures/cross_model_hsweep.pdf}
    \caption{The U-shaped pattern in synthetic experiments. (a) Model accuracy across homophily. MLP (dashed) is constant; GNNs show U-shape. (b) GNN advantage over MLP. GCN exhibits the strongest U-shape (valley: $-18.75\%$ at $h=0.5$), while GraphSAGE shows remarkable stability.}
    \label{fig:ushape_observation}
\end{figure}

\textbf{Key Finding}: GNN performance follows a U-shaped pattern (Figure~\ref{fig:ushape_observation}):
\begin{itemize}
    \item At $h \leq 0.2$ (low homophily): GCN slightly outperforms MLP (+0.3\%)
    \item At $h = 0.5$ (mid homophily): GCN dramatically underperforms MLP (\textbf{-18.75\%})
    \item At $h \geq 0.8$ (high homophily): GCN slightly outperforms MLP (+0.3\%)
\end{itemize}

\begin{table}[t]
\centering
\caption{U-Shape Severity by Model (N=10 runs per setting)}
\label{tab:ushape_severity}
\begin{tabular}{@{}lccc@{}}
\toprule
\textbf{Model} & \textbf{Valley ($h$=0.5)} & \textbf{Peak (extremes)} & \textbf{Amplitude} \\
\midrule
GCN & $-18.75\%$ & $+0.60\%$ & \textbf{19.35\%} \\
GAT & $-5.30\%$ & $+0.55\%$ & 5.85\% \\
GraphSAGE & $-0.10\%$ & $+0.65\%$ & \textbf{0.75\%} \\
\bottomrule
\end{tabular}
\end{table}

\textbf{Cross-Model Analysis} (Table~\ref{tab:ushape_severity}): U-shape severity correlates with aggregation mechanism:
\begin{itemize}
    \item \textbf{GCN} (mean aggregation, replaces features): Strongest U-shape (19.35\% amplitude)
    \item \textbf{GAT} (attention, replaces features): Moderate U-shape (5.85\%)
    \item \textbf{GraphSAGE} (sampling, \textit{concatenates} features): Minimal U-shape (0.75\%)
\end{itemize}

This observation suggests that \textit{feature preservation} (via concatenation) is crucial for robustness. We formalize this insight through the Aggregation Damage Ratio in Section~\ref{sec:adr}.

\subsection{Observation 2: Same Homophily, Opposite Outcomes}
\label{sec:same_h_obs}

If homophily alone determined GNN success, datasets with similar $h$ should show similar GNN advantages. We test this by identifying dataset pairs with nearly identical homophily but different MLP accuracy.

\begin{table}[t]
\centering
\caption{Same Homophily, Different Outcomes: Evidence for Information Budget}
\label{tab:same_h_different_outcome}
\begin{tabular}{@{}llccccc@{}}
\toprule
\textbf{Pair} & \textbf{Dataset} & \textbf{$h$} & \textbf{MLP} & \textbf{Budget} & \textbf{GCN Adv.} \\
\midrule
\multirow{2}{*}{1} & Cora & 0.81 & 75.1\% & 24.9\% & \textbf{+12.7\%} \\
 & Coauthor-CS & 0.81 & 94.4\% & 5.6\% & $-0.7\%$ \\
\midrule
\multirow{2}{*}{2} & PubMed & 0.80 & 87.6\% & 12.4\% & $+0.2\%$ \\
 & Coauthor-Physics & 0.93 & 95.8\% & 4.2\% & $+0.6\%$ \\
\midrule
\multirow{2}{*}{3} & DBLP & 0.83 & 77.8\% & 22.2\% & \textbf{+7.8\%} \\
 & Questions & 0.84 & 97.1\% & 2.9\% & $+0.0\%$ \\
\bottomrule
\end{tabular}
\end{table}

\textbf{Key Finding} (Table~\ref{tab:same_h_different_outcome}): All 7 dataset pairs (full results in Appendix) confirm:
\begin{quote}
\textit{Higher MLP accuracy $\Rightarrow$ Lower GNN advantage, regardless of homophily.}
\end{quote}

The correlation between Budget = $(1 - \text{MLP})$ and GNN advantage is $r = 0.87$ ($p < 0.01$), while the correlation between homophily $h$ and GNN advantage is only $r = 0.31$ ($p = 0.18$) when controlling for MLP accuracy. This directly motivates the Information Budget Principle.

\subsection{Observation 3: The Synthetic-Real Gap}
\label{sec:gap_obs}

The symmetric U-shape in synthetic experiments (Observation 1) raises a question: does it hold in real-world graphs? We conduct semi-synthetic experiments using real node features with rewired graph topology.

\begin{table}[t]
\centering
\caption{Semi-Synthetic Experiments: U-Shape Becomes Monotonic with Real Features}
\label{tab:semi_synthetic_obs}
\begin{tabular}{@{}lcccc@{}}
\toprule
\textbf{Dataset} & \textbf{$h$=0.15} & \textbf{$h$=0.50} & \textbf{$h$=0.75} & \textbf{Pattern} \\
\midrule
\multicolumn{5}{l}{\textit{Synthetic features (controlled separability)}} \\
Synthetic & +6.4\% & $-18.6\%$ & +6.4\% & U-shape \\
\midrule
\multicolumn{5}{l}{\textit{Real features (rewired topology)}} \\
Cora & $-31.4\%$ & $+3.0\%$ & $+24.0\%$ & Monotonic \\
CiteSeer & $-25.2\%$ & $+3.9\%$ & $+10.9\%$ & Monotonic \\
PubMed & $-17.1\%$ & $-8.2\%$ & $+5.1\%$ & Monotonic \\
\bottomrule
\end{tabular}
\end{table}

\textbf{Key Finding} (Table~\ref{tab:semi_synthetic_obs}): With real features:
\begin{itemize}
    \item \textbf{No U-shape}: GNN advantage increases monotonically with $h$
    \item \textbf{Crossover at $h \approx 0.5$}: GNN transitions from losing to winning
    \item \textbf{Low-$h$ is harmful}: Unlike synthetic data, real low-$h$ graphs hurt GNN
\end{itemize}

\textbf{Why the Gap?} In synthetic data, low-$h$ means neighbors have \textit{opposite} features (exploitable contrastive signal). In real data, low-$h$ neighbors have \textit{orthogonal} features (noise). We formalize this distinction as Dual Heterophily Types in Section~\ref{sec:heterophily}.

\subsection{Summary: Three Puzzles Requiring Explanation}

These observations reveal three puzzles that homophily alone cannot explain:

\begin{enumerate}
    \item \textbf{U-Shape Puzzle}: Why does GNN performance form a U-shape in synthetic data?
    \item \textbf{Same-$h$ Puzzle}: Why do same-homophily datasets show opposite GNN outcomes?
    \item \textbf{Synthetic-Real Puzzle}: Why does U-shape become monotonic with real features?
\end{enumerate}

The Information Budget framework (Section~\ref{sec:budget}--\ref{sec:heterophily}) provides unified answers:
\begin{itemize}
    \item \textbf{Puzzle 1}: U-shape reflects aggregation damage at mid-$h$, quantified by ADR
    \item \textbf{Puzzle 2}: MLP accuracy determines the budget ceiling for GNN improvement
    \item \textbf{Puzzle 3}: Real features create Type B heterophily (irrecoverable), unlike synthetic Type A
\end{itemize}

Having established these empirical patterns, we now develop the rigorous SNR framework (Section~\ref{sec:snr_theory_rigorous}) that explains them through the lens of discriminability dynamics.
